\documentclass[10pt,a4paper]{article}
\usepackage[utf8]{inputenc}
\usepackage[T1]{fontenc}
\usepackage[french]{babel}
\usepackage[a4paper,margin=1.8cm]{geometry}
\usepackage{booktabs,array,float,hyperref}
\hypersetup{colorlinks=true,linkcolor=blue,urlcolor=blue}
\emergencystretch=1em

\title{Pipeline causal des données réelles\\Phase 3/10}
\author{Stéganographie adaptative par parcours de graphe}
\date{13 juin 2026}

\begin{document}
\maketitle

\begin{abstract}
Cette phase construit le socle de données réelles de l'étude. Cinq datasets
sont acquis, vérifiés par checksum et convertis en Parquet : GeoLife, T-Drive,
tgbl-wiki, MOOC et LastFM. Deux schémas canoniques séparent trajectoires et
interactions temporelles. Un validateur indépendant confirme l'ordre causal
des splits et l'absence de séquence commune entre entraînement, validation et
test. Les statistiques révèlent des différences fortes d'entropie
conditionnelle entre domaines, ce qui renforce la nécessité d'un encodage
adaptatif avec abstention.
\end{abstract}

\section{Panel réel}

\begin{table}[H]
\centering
\caption{Datasets effectivement traités.}
\begin{tabular}{lllr}
\toprule
Dataset & Niveau & Domaine & Taille utile\\
\midrule
GeoLife 1.3 & A & mobilité personnelle & 24 867 361 points\\
T-Drive & A & taxi urbain & 10 055 095 points\\
tgbl-wiki-v2 & B & édition collaborative & 157 474 événements\\
MOOC/JODIE & B & apprentissage en ligne & 411 749 événements\\
LastFM/JODIE & B & écoute musicale & 1 293 103 événements\\
\bottomrule
\end{tabular}
\end{table}

Le niveau A contient des séquences observées et porte les conclusions
principales de naturalité. Le niveau B sert à l'apprentissage temporel, au
passage à l'échelle et au transfert; il ne permet pas d'affirmer qu'un parcours
reconstruit est un comportement humain naturel.

\section{Schémas et découpage}

Le schéma d'interaction contient les identifiants de l'événement, de la source,
de la destination et de la séquence, puis le timestamp, le label et le split.
Le schéma de trajectoire contient les identifiants du point, de la séquence et
de l'entité, puis le timestamp, les coordonnées et le split. Les utilisateurs
et items reçoivent des préfixes distincts.

Les flux TGB/JODIE suivent un découpage 70/15/15 par timestamp, en conservant
ensemble les égalités. GeoLife est partagé par trajectoire complète. Deux
trajectoires traversant une frontière, soit 9 616 points, sont mises en
quarantaine. T-Drive est segmenté aux changements de jour et après 20 minutes
d'interruption. Ses frontières sont placées à minuit : aucune session ne
traverse un split.

Le validateur relit les Parquet, recalcule les checksums, compare les ensembles
d'identifiants et impose
\[
 \max t_{\rm train}<\min t_{\rm val}
 \leq\max t_{\rm val}<\min t_{\rm test}.
\]
Cette propriété est vérifiée pour GeoLife et T-Drive.

\section{Structure temporelle}

\begin{table}[H]
\centering
\caption{Statistiques des flux bipartites.}
\begin{tabular}{lrrrr}
\toprule
Dataset & Événements & Noeuds & Arêtes répétées & $H(D\mid S)$\\
\midrule
tgbl-wiki & 157 474 & 9 227 & 88,4\% & 1,197\\
MOOC & 411 749 & 7 144 & 56,7\% & 4,880\\
LastFM & 1 293 103 & 1 980 & 88,0\% & 7,034\\
\bottomrule
\end{tabular}
\end{table}

Pour tgbl-wiki, \(2^{H(D\mid S)}=2,293\) continuations effectives en moyenne.
LastFM offre une incertitude conditionnelle bien supérieure. L'avis est
tranché : un débit fixe ou un seuil d'entropie unique ne sera pas défendable
sur ce panel. Le contrôleur devra adapter localement sa décision et accepter
l'abstention dans les contextes concentrés.

GeoLife contient 18 668 trajectoires retenues, avec une médiane de 506 points,
une durée médiane de 2 527,5 secondes et un intervalle médian de deux secondes.
T-Drive produit 211 404 sessions, avec 20 points, 5 408 secondes et 300
secondes respectivement. Ces dynamiques sont suffisamment différentes pour
tester la robustesse de la modélisation temporelle.

\section{Qualité et provenance}

GeoLife ne contient qu'un point invalide et aucune inversion temporelle. Son
archive comporte 18 670 fichiers PLT alors que le guide en annonce 17 621;
l'écart est documenté. T-Drive contient 98 lignes invalides et 33 158
coordonnées hors de l'enveloppe large Chine
\(\mathrm{lat}\in[18,54]\), \(\mathrm{lon}\in[73,135]\), toutes exclues.

L'échantillon T-Drive officiel disponible présente 8 911 fichiers et 10,09
millions de points bruts, contre 10 357 taxis et environ 15 millions de points
dans l'article. Il ne sera donc pas décrit comme la collection complète.

\section{Licences et limites}

GeoLife autorise uniquement l'usage non commercial. La page T-Drive ne donne
pas de termes explicites de redistribution. Les conditions SNAP de MOOC et
LastFM doivent être validées avant publication du dépôt. Porto Taxi reste
conditionné par l'authentification Kaggle.

La discrétisation spatiale est volontairement différée. Les cellules, le
vocabulaire et les normaliseurs devront être ajustés uniquement sur
l'entraînement; les choisir sur l'ensemble des points introduirait une fuite.

\section{Décision de passage}

La phase 3 est techniquement validée : cinq datasets réels, quatre familles
comportementales, 16 tests réussis, fiches de provenance, checksums, sorties
Parquet et validation causale indépendante. La phase 4 peut commencer avec
tgbl-wiki pour le développement, GeoLife et T-Drive pour les parcours
observés, puis MOOC et LastFM pour le transfert.

Le paquet public Q1 restera conditionné à une décision de redistribution pour
T-Drive et SNAP, ou à leur remplacement par des sources aux termes explicites.

\end{document}
